% Chapter 4, typeset from lectures/L04: the README is the opener and the closing review, and each
% part of Appendix A is a section, in the same order and with the same subsections. Appendix B,
% the exercises, is the last section.

\chapter{Metastability and Synchronization}
\label{c4:ch}

{\large\itshape Making an asynchronous input safe to use in a synchronous system.\par}

\begin{chapterbox}{In this chapter}
\begin{itemize}
  \item Why an asynchronous input can break the clean edge detectors and registers built in
    Chapter~\ref{c3:ch}.
  \item Metastability: what it is, and why it cannot simply be designed away.
  \item The double-flop synchronizer, and the assert-asynchronously, release-synchronously reset
    pattern.
  \item Reusing the same chain for edge detection, and what it does and does not do about bounce.
  \item Generics: one module serving several sizes.
  \item \code{led_toggle_sync}, and what it does on the DE0-CV when a real, bouncing push button
    drives it.
  \item Setup/hold timing detail (\secref{c4:a4}), and what a real debounce adds (\secref{c4:a7}).
\end{itemize}
\end{chapterbox}

This chapter takes Chapter~\ref{c3:ch}'s LED toggle and makes it safe for a real button. The course
builds it live, in this order: first as a gate network in CircuitVerse, with two flip-flops in
front of the button, two in front of the reset, and the edge detector behind them; then in VHDL,
as \code{reset_sync}, which asserts asynchronously and releases synchronously, and
\code{button_sync}, the three-flop chain with a \code{generic} so that one file serves a design
with one button and a design with two; then \code{led_toggle_sync}, the two subcomponents wired
together around the toggle logic, which is all that is left of the original circuit; and finally
onto the board. Before the synchronizer goes in, predict what a real, asynchronous button does to
Chapter~\ref{c3:ch}'s edge detector. That prediction is the reason for the whole chapter.

The chapter opens with metastability itself, and the argument worth holding on to is why two
flip-flops make the failure \emph{unlikely} rather than impossible: that is the part that does not
survive being remembered as a rule of thumb. Two parts of the material are ones the course leaves
to reading rather than presenting: the setup/hold timing detail in \secref{c4:a4}, beyond the fact
that the window exists and that violating it is what causes the trouble; and \secref{c4:a7}, on
what a real debounce circuit adds. Both come back in the review questions and are drilled by
Exercises~\ref{c4:ex:2} and~\ref{c4:ex:7}, so they are not optional.

Before you read on, be comfortable with D flip-flops, clocking and edge detection from
Chapter~\ref{c3:ch}; this chapter assumes all of it. It also builds on \secref{c2:a6}: its worked
example is composed from two subcomponents, and \secref{c4:a8} adds the \code{generic} that lets one
of them serve designs with different numbers of buttons. Afterwards, work through the exercises and
read \file{lectures/L04/led_toggle_sync/led_toggle_sync.vhd}. Its \code{reset_sync} and
\code{button_sync} subcomponents are Exercises~\ref{c4:ex:8} and~\ref{c4:ex:9} rather than files in
the repository, so do those first, copy your own into \file{led_toggle_sync/}, and then run its
testbench. If this chapter does not click on the first read, Erik Pihl's video walkthrough of
metastability and the double-flop synchronizer (\url{https://www.youtube.com/watch?v=KrssJRgF13I})
and the metastability pages at nandland (\url{https://nandland.com/lesson-13-metastability/}) and
VHDLwhiz (\url{https://vhdlwhiz.com/terminology/metastability/}) cover the same concept from a
slightly different angle.

% Sections 4.1 to 4.9, from lectures/L04/appendix/a_metastability_and_synchronization.md: one
% section per part, A.1 to A.9.

\section{The Problem: A Button Isn't a Testbench Signal}
\label{c4:a1}
In Chapter~\ref{c3:ch} you built edge detectors and reasoned about them against a clean clock, where
every input changed neatly between edges, exactly the way a simulator draws it. A push button, an
external reset line, or any signal from another clock domain changes whenever the outside world
decides to, completely independent of your FPGA's clock.

Every D flip-flop has two timing requirements around the active clock edge:
\begin{itemize}
  \item a \textbf{setup time}: how long \code{D} must be stable \emph{before} the edge.
  \item a \textbf{hold time}: how long \code{D} must stay stable \emph{after} it.
\end{itemize}
If an input changes inside that window the flip-flop is not guaranteed to capture a clean \code{0}
or \code{1}, and you can never promise it will land outside, because you do not control when an
external signal changes. The rest of this chapter solves that not by avoiding the problem but by
containing it.

\section{Metastability: What Actually Happens Inside the Flip-Flop}
\label{c4:a2}
A D flip-flop is built from cross-coupled gates that snap to \code{0} or \code{1} the instant the
clock samples the input. If \code{D} changes too close to the edge, the internal feedback loop can
be left balanced between the two states, neither fully high nor fully low, for longer than
expected. That is \textbf{metastability}.

Picture a ball balanced exactly on top of a hill. Given enough time it rolls down one side, but
exactly when it commits, and to which side, is not predictable from the moment it was placed there,
and while it is still balanced its position is not a valid digital value at all.

A metastable output behaves the same way:
\begin{itemize}
  \item it almost always resolves to a valid \code{0} or \code{1} on its own, given enough time.
  \item ``enough time'' is not free, and is not the same on every clock cycle.
  \item if the unresolved voltage reaches downstream logic before it settles, different gates can
    legitimately read it differently: one sees \code{0}, another \code{1}, at the same instant.
\end{itemize}
That disagreement is where real systems misbehave: a counter that increments in one part of the
circuit but not another, or a state machine in an impossible state. Everything in this chapter
exists to make sure metastability resolves \emph{before} a signal touches the rest of your design.

In time, the whole problem looks like this:

\bookfigure{setup_hold}{\code{D} changes inside the setup/hold window around a rising clock edge,
leaving \code{Q} metastable before it settles, late, to either level.}{c4:fig:setup_hold}

\noindent Read it left to right. The shaded band is the window from \secref{c4:a1}, \code{t_su}
before the edge and \code{t_h} after it. \code{D} changes inside it, which is exactly what an input
you do not control is free to do. \code{Q} is then neither \code{0} nor \code{1} for a while, and
when it does settle, \emph{when} it settles and \emph{which level} it settles to are both outside
your control. Both dashed continuations are legitimate outcomes of the same input.

Note what the diagram does not show: a wrong answer. Sampling a signal mid-transition is allowed to
give you either level, and that on its own is harmless, because the input genuinely was changing at
that instant. The problem is the delay before you get an answer at all, and that is what the next
section buys time for.

\section{The Double-Flop Synchronizer}
\label{c4:a3}
The standard fix is the \textbf{double-flop synchronizer}:
\begin{itemize}
  \item Place two D flip-flops in series on every asynchronous input: every input that changes with
    no fixed relationship to your clock. A push button, a reset line, a signal arriving from
    another clock domain.
  \item An input that is already clocked by \emph{this} clock is not asynchronous, and must
    \textbf{not} be synchronized. Doing so only delays it, and shifts its timing relative to
    everything around it. \code{seq_detect_mealy} in Chapter~\ref{c8:ch} is exactly that case: its
    \code{din} is synchronous serial data, so it goes straight into the machine while
    \code{reset_n} still gets a synchronizer.
  \item Only ever use the output of the \emph{second} flip-flop elsewhere in the design.
\end{itemize}

\begin{textdrawing}
async_in
   │
   ▼
 [FF1] ──────► [FF2] ──────► synchronized, safe to use
   (may go        (extremely likely
   metastable)     to be stable)
\end{textdrawing}

\noindent The same two flip-flops, in time:

\bookfigure{synchronizer_waveform}{\code{async_in} changes just before a clock edge; \code{s1} may go
metastable for a clock period, and \code{s2} comes out clean one edge later.}{c4:fig:synchronizer}

\noindent\code{FF1} is the one exposed to the asynchronous input, so it is the one that might go
metastable. Its output only has to reach \code{FF2}, and it has an entire clock period to get there.
The probability that a metastable node is still unresolved after a full period is small, and
shrinks the longer it is given, so by the time \code{FF2} samples on the next edge, \code{FF1} has
almost always settled. \code{FF2}'s output is what you build the rest of your synchronous logic on.

This is probabilistic, not absolute, and \secref{c4:a4} says what that means. At the speeds used in
this book, tens of MHz, two stages are the standard, industry-accepted answer.

One rule ties it together: \textbf{every asynchronous input gets its own double-flop chain before
it is used anywhere else.} External reset buttons, push buttons, and later any
\textbf{single-bit} signal crossing from another clock domain.

That last word is load-bearing, and it is the one limit worth learning before you over-apply the
rule. A double-flop chain synchronizes \textbf{one bit}. Put a multi-bit bus through a chain per bit
and each bit is individually safe but they are not safe \emph{together}: each independently
resolves to the old or new value on the same edge, so the receiver can sample a combination that
was never on the bus. A counter stepping from \code{0111} to \code{1000} has every bit changing at
once, and a per-bit synchronizer can hand you \code{1111} or \code{0000}, values the counter never
held.

Buses crossing clock domains need a handshake, a FIFO, or Gray coding, which arranges that only one
bit ever changes at a time. This book never crosses clock domains with a bus, so you will not need
them here. Knowing that the two-flop answer stops at one bit is enough for now.

\section{Timing Intuition: Why ``Almost Certainly'' Is Good Enough}
\label{c4:a4}
A synchronizer does not buy ``no metastability''. It buys metastability that gets a full clock
period to resolve before it can affect anything else, and that is a statistical argument rather
than a proof. The correct posture is ``reasonably certain, backed by an extremely low,
well-characterized failure rate'', not ``provably safe under all conditions'', which is why extra
stages exist for extreme cases instead of two being declared sufficient by law of physics.

The relationship has a standard shape, worth seeing once even though you will not compute with it
here:

\begin{console}
MTBF = e^(t_r / tau) / (T_w * f_clock * f_data)
\end{console}

\noindent Read it for what it says rather than for numbers. \code{t_r} is the resolution time a
stage gets, roughly one clock period; \code{tau} and \code{T_w} are both properties of the
flip-flop, \code{tau} being how fast a metastable value decays and \code{T_w} the width of the
window around the clock edge in which the input can provoke one at all. Adding a stage buys another
\code{t_r} in the exponent, so \textbf{the numerator grows exponentially with every stage you add},
against one clock cycle of added latency. The denominator grows linearly with the clock rate and,
just as importantly, with \code{f_data}: how often the asynchronous input actually changes. A
button pressed once a second and a data line toggling at 10\,MHz differ by seven orders of magnitude
in that term alone, which is why the case for a third stage is a question about the input, not
about the clock.

\section{Synchronizing a Reset Signal: Assert Async, Release Sync}
\label{c4:a5}
Resets deserve a specific pattern, which \code{reset_sync} demonstrates. There is no
\file{reset_sync.vhd} in the repository to open: it is Exercise~\ref{c4:ex:8}, and you write it.
The listing below is the whole of its logic; the \code{library}/\code{use} clauses, the
\code{entity}, the \code{architecture} frame and the declaration of the internal \code{reset_s1_n}
are yours to add around it.

\begin{vhdlcode}
process(clock, reset_n) is
begin
    if (reset_n = '0') then
        reset_s1_n <= '0';
        reset_s2_n <= '0';
    elsif (rising_edge(clock)) then
        reset_s1_n <= '1';
        reset_s2_n <= reset_s1_n;
    end if;
end process;
\end{vhdlcode}

\begin{itemize}
  \item \textbf{Assert asynchronously.} When \code{reset_n} drops, both stages are forced to
    \code{0} immediately, with no dependency on the clock. A reset that waited for a clock edge
    would defeat its own purpose: it must take effect the instant it is needed, clock or no clock.
  \item \textbf{Release synchronously.} When \code{reset_n} returns to \code{1} the reset condition
    does not vanish everywhere at once. It shifts out through the two stages like any other
    asynchronous input, so \code{reset_s2_n}, the signal the design actually uses, releases cleanly
    on a clock edge together with everything else. Without this, different flip-flops could leave
    reset on different cycles and momentarily disagree about system state, the same class of bug
    metastability causes.
\end{itemize}
This is why \code{reset_n} is synchronized by its own chain before anything else depends on it,
including the button synchronizer, whose reset input is \code{reset_s2_n} rather than the raw
\code{reset_n}.

\section{Reusing the Chain for Edge Detection --- and, Incidentally, Debouncing}
\label{c4:a6}
The button needs synchronizing, and from Chapter~\ref{c3:ch} you also need to detect \emph{when} it
is pressed rather than its level. \code{button_sync}, which is Exercise~\ref{c4:ex:9} and so is also
yours to write rather than a file in the repository, gets both from a single chain of
\textbf{three} flip-flops:

\begin{textdrawing}
button_n
   │
   ▼
[FF1] → [FF2] → [FF3]
        │        │
        │        └─ previous value
        └─ current value
\end{textdrawing}

\noindent Flip-flops 1 and 2 are \secref{c4:a3}'s synchronizer, making \code{button_s2_n} the stable
current value. Flip-flop 3 stores last cycle's \code{button_s2_n} as \code{button_s3_n}, the
previous value. Flip-flops 1 and 2 are worth noticing on their own: Chapter~\ref{c5:ch}'s
Exercise~\ref{c5:ex:4} extracts that double-flop chain into a reusable \code{sync} module, leaving
the third flip-flop here as the edge detector's own. Comparing stage 2 against stage 3 gives edge
detection for free:

\begin{center}
\code{button_edge_s2 = (not button_s2_n) and button_s3_n}
\end{center}

\noindent which is true for exactly one clock cycle: the one where the button reads pressed now and
did not the cycle before.

\textbf{Synchronization, on its own, only solves metastability.} It says nothing about a mechanical
button's tendency to bounce: to make and break contact several times over the first millisecond or
so, each bounce looking like a separate press. What keeps this book's exercises from being swamped
by bounce is the \emph{combination} of the chain with a detection window that is slow relative to
the bounce:
\begin{itemize}
  \item In this chapter's CircuitVerse exercises the clock period is a full second. Bounce settles
    in far less than that, so by the next edge \code{button_s2_n} reads one clean level. The chain
    looks like it is debouncing, but it is really just sampling too infrequently to see bounce.
  \item On the FPGA the clock is \code{50 MHz}, one sample every \code{20 ns}, while bounce lasts
    from under a millisecond to a few tens of milliseconds, many thousands of cycles. At that speed
    \code{button_sync} faithfully detects \emph{every} bounce transition as its own edge, and the LED
    could toggle several times for one physical press.
\end{itemize}
So precisely: the double-flop synchronizer reliably solves metastability at the speeds used here.
Whether it also \emph{looks} like debouncing depends entirely on the clock period being slow
relative to the switch's bounce time, which is true in CircuitVerse and not automatically true on
the FPGA.

\section{What Real Hardware Debouncing Adds}
\label{c4:a7}
This book has no dedicated debounce chapter, and the worked examples deliberately keep using
\secref{c4:a6}'s circuit, bounce and all. It is good enough for a demo LED. A production design adds
one of the following \emph{on top}:
\begin{itemize}
  \item \textbf{An analog RC low-pass filter}, optionally with a Schmitt-trigger buffer, on the
    button input before it reaches an FPGA pin. The RC pair smooths the fast make/break transitions
    into a single gradual edge, and the Schmitt trigger, with its two switching thresholds, turns
    that back into one clean digital transition. This solves bounce before the signal is even
    digital, but costs board components and does not help for inputs that are not mechanical
    switches.
  \item \textbf{A digital debounce counter.} Rather than acting on the next detected edge, require
    the input to hold a new level for some minimum duration, a few milliseconds' worth of clock
    cycles, before accepting it. That builds your own artificially slow detection window on-chip,
    the same effect the 1-second CircuitVerse clock gives for free, without slowing the system
    clock. You will meet counters in Chapter~\ref{c6:ch} and timers, which this technique needs, in
    Chapter~\ref{c7:ch}.
\end{itemize}
Both sit on top of the synchronizer, not instead of it: even a perfectly debounced signal is still
asynchronous to your clock and still needs a double-flop chain.

\section{Generics: One Module, Several Sizes}
\label{c4:a8}
\code{reset_sync} and \code{button_sync} are built from the same idea, but only one has a question
to answer. A reset is one bit and always will be. How many \emph{buttons} does a design have? This
chapter's \code{led_toggle_sync} has one, Chapter~\ref{c8:ch}'s \code{fsm_led} has two. Writing the
same three-flop chain out once per width is exactly the duplication that splitting them into
modules set out to avoid.

A \textbf{generic} answers it: a value an entity declares alongside its ports, fixed when the design
is built rather than changing while it runs.

\begin{vhdlcode}
entity button_sync is
    generic(COUNT: natural range 1 to 3 := 1);
    port(clock, reset_s2_n: in  std_logic;
         button_n         : in  std_logic_vector(COUNT-1 downto 0);
         button_edge_s2   : out std_logic_vector(COUNT-1 downto 0));
end entity;
\end{vhdlcode}

\begin{itemize}
  \item \code{COUNT} behaves like a constant inside the architecture. It cannot change while the
    circuit runs; it is decided before a single gate is placed.
  \item \code{:= 1} is its default, used by any instantiation that says nothing.
  \item The port widths are written \emph{in terms of it}, and so is every internal signal that
    follows.
\end{itemize}

\subsection{Overriding it}
\label{c4:override}
An instantiation passes generics with a \code{generic map}, before the \code{port map}:

\begin{vhdlcode}
button_sync1: entity work.button_sync
    generic map(2)
    port map(clock, reset_s2_n, button_n, button_edge_s2);
\end{vhdlcode}

\noindent Generics are matched \textbf{by position}, in declaration order, exactly as ports are in
\secref{c2:a6}. Leaving the \code{generic map} out applies the default; this book writes it
explicitly even when it matches, because a reader should not have to open the module to find out
how wide an instance is.

\subsection{When a generic is worth having}
\label{c4:genericworth}
\code{button_sync} has one and \code{reset_sync} has none, and that difference is the whole rule:
\textbf{add a generic when something genuinely varies between instances, not on principle.}

\code{button_sync} is instantiated at \code{COUNT = 1} by \secref{c4:a9}'s \code{led_toggle_sync}
and at \code{COUNT = 2} by Chapter~\ref{c8:ch}'s \code{fsm_led}: two widths, one file.
\code{reset_sync} is one bit everywhere, so a \code{WIDTH} generic on it would be a parameter nobody
ever passes anything but \code{1} to, and one more thing for a reader to check. That second
instantiation is what makes a generic worth having; a parameter with exactly one value in the
entire book would just be indirection. \code{button_sync}'s own testbench instantiates it at
\emph{both} widths in one simulation, precisely to prove that one file really does serve both.

You will meet one more generic in Chapter~\ref{c7:ch}: \code{timer}'s \code{TICK_COUNT}, which is a
count rather than a width, and is what lets the same file measure a fiftieth of a second in one
design and a whole second in another.

\section{Worked Example: \code{led_toggle_sync}}
\label{c4:a9}
This ties \secref{c4:a1} to \secref{c4:a6} together into Chapter~\ref{c3:ch}'s toggle circuit, made
safe for a real, asynchronous button and reset.

\begin{booktable}{@{}p{6em}p{5em}L@{}}
  \thead{Port} & \thead{Direction} & \thead{Description} \\
  \midrule
  \code{clock} & in & \code{50 MHz} system clock on the DE0-CV board. \\
  \code{reset_n} & in & Active-low asynchronous reset, wired to a push button. \\
  \code{button_n} & in & Active-low push button; toggles the LED on its falling edge. \\
  \code{led} & out & LED, toggled on each synchronized, edge-detected button press. \\
\end{booktable}

\noindent Internal signals: \code{led_s} holds the LED state, \code{reset_s2_n} is the reset after
the two-flop synchronizer (the \code{s2} suffix is this chapter's convention for ``synchronized
with two flip-flops''), and \code{button_edge_s2} is the one-cycle pulse marking a synchronized
falling edge on \code{button_n}.

Like Chapter~\ref{c2:ch}'s two-digit display, this design is built from more than one entity, using
the instantiation from \secref{c2:a6} and the generic from \secref{c4:a8}. The difference is what
the pieces do: the display instantiated one module twice to do the same job on two halves of its
input, while this instantiates two \emph{different} modules with one job each. Both are the same
mechanism. Composing a design out of modules is not a special technique for repeated logic; it is
how designs are built once they outgrow a single file.

The two chains live in two subcomponents:
\begin{itemize}
  \item \code{reset_sync} is \secref{c4:a5}'s two-flop reset synchronizer. It takes the raw
    \code{reset_n} and returns \code{reset_s2_n}.
  \item \code{button_sync} is \secref{c4:a6}'s three-flop synchronizer-plus-edge-detector. It takes
    the already synchronized \code{reset_s2_n} and pulses \code{button_edge_s2} once per press.
\end{itemize}
Keeping them apart matters more than it looks. A design with no button of its own still needs its
reset synchronized, and with two modules it simply instantiates the one it wants, as
Chapter~\ref{c8:ch}'s \code{seq_detect_mealy} does. Welded into one module, that design would have
to instantiate the whole thing and disable the half it does not use.

\code{button_sync} takes a \emph{vector} of buttons so one file can serve a two-button design. This
design has one, so \file{led_toggle_sync.vhd} declares a pair of one-element vectors to connect to
it:

\begin{vhdlcode}
signal button_n_v, button_edge_s2_v: std_logic_vector(0 downto 0);
...
button_n_v(0)  <= button_n;
button_edge_s2 <= button_edge_s2_v(0);
\end{vhdlcode}

\noindent That is plumbing, not logic: two wires renamed.

The top-level architecture contains nothing but the toggle process itself, and it only ever touches
the already-synchronized \code{reset_s2_n} and \code{button_edge_s2}, never the raw inputs:

\begin{vhdlcode}
LED_PROCESS: process(clock, reset_s2_n) is
begin
    if (reset_s2_n = '0') then
        led_s <= '0';
    elsif (rising_edge(clock)) then
        if (button_edge_s2 = '1') then
            led_s <= not led_s;
        end if;
    end if;
end process;
\end{vhdlcode}

\noindent The full source of the top level is \file{lectures/L04/led_toggle_sync/led_toggle_sync.vhd}:

\vhdlfile{lectures/L04/led_toggle_sync/led_toggle_sync.vhd}

\noindent The two subcomponents are not in the repository, because writing them is
Exercises~\ref{c4:ex:8} and~\ref{c4:ex:9}. What this chapter gives you of each is:
\code{reset_sync}'s process in \secref{c4:a5} above, and \code{button_sync}'s chain as a diagram and
an equation in \secref{c4:a6} with its entity in \secref{c4:a8}. The surrounding entity and
architecture in the first case, and the whole architecture body in the second, are the exercises'
own work. Copy your own into \file{led_toggle_sync/} to build it.

\textbf{By hand, in CircuitVerse first.} Before touching VHDL, build it as a gate network: two D
flip-flops in series for the reset synchronizer; two for \code{button_n} plus a third storing the
previous synchronized value; an AND gate combining the inverted current value with the previous one
to produce \code{button_edge_s2}; and a toggle flip-flop for \code{led_s}, its inverted output fed
back to its input, enabled only when \code{button_edge_s2 = 1}. Set the clock period to
\mbox{\code{1000 ms}} so you can watch each stage by eye.

\bookimage{1.0}{lectures/L04/appendix/images/led_toggle_sync_circuit.png}{The synchronized LED
toggle in CircuitVerse: the reset and button double-flop synchronizers, the edge-detecting AND gate,
and a toggle flip-flop for \code{led_s}.}{c4:fig:circuit}

\noindent One thing in the drawing needs reading carefully: the flip-flops carry a \code{preset}
input, and on the three button-synchronizer stages it is tied to \code{1}. That is not decoration.
In CircuitVerse, \code{preset} is the value a flip-flop loads when its reset is asserted, so tying
it to \code{1} is what makes \code{button_s1_n}, \code{button_s2_n} and \code{button_s3_n} come out
of reset holding \code{1} rather than \code{0}.

That matches \code{button_sync}, which resets the same three stages to \code{(others => '1')}, and
it matters for the same reason: the buttons are active-low, so \code{1} means ``released''. Reset
the chain to \code{0} instead and it comes up claiming the button is already held down. The reset
flip-flops leave \code{preset} alone, because there a reset to \code{0} really is all the circuit
asks for.

Press the button in the simulation and watch it propagate. The LED does not toggle instantly: it
waits for the synchronized, edge-detected pulse to reach the toggle flip-flop, which takes a small,
fixed number of clock pulses, exactly the latency \secref{c4:a3} called the price of safety.

\textbf{On the FPGA:} the course demonstrates the same design on the Terasic DE0-CV (device
\code{5CEBA4F23C7N}), with \code{clock} on the \code{50 MHz} oscillator, \code{reset_n} and
\code{button_n} on two push buttons, and \code{led} on an LED; Appendix~\ref{app:quartus} is the
workflow, for a reader who has a board. The LED toggles on every press, and per \secref{c4:a6} a
real press may occasionally produce more than one toggle from contact bounce.

% Section 4.10, from lectures/L04/README.md: the objectives, the questions to test yourself, and the
% next lecture.

\needspace{10\baselineskip}
\section{Review}
\label{c4:sec:review}

\subsection*{What you should be able to do now}
After this chapter, you should be able to:
\begin{itemize}
  \item Explain what metastability is, why an asynchronous input can cause it, and why two
    flip-flops make it unlikely rather than impossible.
  \item Apply the double-flop synchronizer to any asynchronous input, including the ``assert
    asynchronously, release synchronously'' pattern used for resets.
  \item Say why that chain doubles as a practical but incomplete debounce, and what a real one would
    add.
  \item Give a module a \code{generic}, override it with a \code{generic map}, write a port's width
    in terms of it, and judge when a generic earns its place rather than adding one on principle.
\end{itemize}

\needspace{6\baselineskip}
\subsection*{Questions to test yourself}
You should be able to explain:
\begin{enumerate}
  \item What exactly happens when an asynchronous input to a flip-flop changes too close to an
    active clock edge?
  \item How does the double-flop synchronizer address that, and why is the \emph{second}
    flip-flop's output the one that is safe to use?
  \item Why does the chain used here behave like a debounce circuit in CircuitVerse but not fully
    debounce a real button at \code{50 MHz}, and why is that a property of the clock period rather
    than of the circuit? What would you add to fix it?
  \item \code{button_sync} has a generic and \code{reset_sync} has none. What is the rule, and what
    would be wrong with giving \code{reset_sync} a \code{WIDTH} generic ``for symmetry''?
\end{enumerate}

\needspace{6\baselineskip}
\subsection*{Looking ahead}
Chapter~\ref{c5:ch} is about the one construct left, and what the toolchain builds:
\begin{itemize}
  \item \code{variable}: the one construct you have not met yet, and the sharp distinction between
    how it updates and how a \code{signal} does.
  \item Why the same accumulation written with a \code{signal} and with a \code{variable} gives two
    different answers, and which of the two is the bug.
  \item What the toolchain actually builds: LUTs and flip-flops, what limits a clock's speed, and
    the latch you infer by accident.
\end{itemize}

% Section 4.11, from lectures/L04/appendix/b_exercises.md.

\section{Exercises}
\label{c4:sec:exercises}

\begin{aside}
\textbf{How to work these.} Each exercise is tagged with the kinds of work it asks for:
\kinds{Paper} by hand, \kinds{CircuitVerse} in the browser, \kinds{VHDL} in a module of your own.

\textbf{How to check your work.} Every exercise that asks you to write a VHDL module ships a
self-checking testbench under \file{lectures/L04/exercises} in the companion repository. Write your
module in its \file{exercises/<module>} directory, using the entity name and the \textbf{port
order} the exercise specifies, then run it with GHDL; \secref{c2:app:c} has the three commands.

No FPGA board is needed for any exercise. Your job is to get the VHDL right, and the testbench is
how you confirm it. Short answers to the parts you work on paper are in
Appendix~\ref{app:answers}; reference solutions to the VHDL are published in the companion
repository, under \file{lectures/L04/solutions}. Try each exercise yourself before looking at
either.
\end{aside}

% ------------------------------------------------------------------------------------------
\exerciseset{Metastability Concepts}
\label{c4:set:1}

\exercise{Metastability in your own words}{Paper}
\label{c4:ex:1}
Explain metastability in your own words. In your explanation, describe:
\begin{itemize}
  \item What happens inside a D flip-flop when its input changes too close to the active clock
    edge.
  \item Why the flip-flop output may temporarily remain between a valid logic \code{0} and logic
    \code{1}.
  \item Why the output may take an unpredictable amount of time to settle.
  \item Why two downstream gates may interpret the unresolved output differently:
    \begin{itemize}
      \item One gate may interpret it as \code{0}.
      \item Another gate may interpret it as \code{1}.
    \end{itemize}
\end{itemize}
Explain why this disagreement can cause different parts of a digital circuit to enter inconsistent
states.

\exercise{The setup-and-hold window}{Paper}
\label{c4:ex:2}
Setup time and hold time define a window around the active clock edge during which the input of a
flip-flop must remain stable.

\noindent Explain:
\begin{itemize}
  \item Why a synchronously generated signal can normally be designed to satisfy these timing
    requirements.
  \item Why an asynchronous input, such as a mechanical pushbutton, has no fixed relationship to the
    system clock.
  \item Why an asynchronous input may therefore change during the setup-and-hold window regardless
    of how the surrounding synchronous logic is designed.
\end{itemize}
Conclude why the metastability risk from an asynchronous input must be managed rather than simply
avoided.

% ------------------------------------------------------------------------------------------
\exerciseset{Synchronizer Design}
\label{c4:set:2}

\exercise{Drawing a double-flop synchronizer}{Paper}
\label{c4:ex:3}
Draw a double-flop synchronizer for one asynchronous input.

\noindent The circuit must contain:
\begin{itemize}
  \item One asynchronous input.
  \item Two D flip-flops connected in series.
  \item A shared system clock.
  \item A synchronized output taken from the second flip-flop.
\end{itemize}
Explain in your own words:
\begin{itemize}
  \item Why the first flip-flop is exposed directly to the asynchronous input.
  \item Why the first flip-flop may become metastable.
  \item How the interval between the first and second clock edges gives the first stage time to
    settle.
  \item Why the output of the second flip-flop has a much lower probability of being metastable.
  \item Why the second stage, rather than the first, should be used by the rest of the design.
\end{itemize}
A synchronizer does not make the probability of metastability exactly zero. Explain why it
nevertheless reduces the risk to an acceptably low level in most designs.

\note[Tip] Consider how much resolution time the first flip-flop receives before its value is
sampled by the second flip-flop.

\exercise{A third synchronizer stage}{Paper}
\label{c4:ex:4}
A colleague suggests using three synchronizer flip-flops on every asynchronous input in a low-speed
design clocked at a few megahertz, just to be safe.

A third stage buys one more clock period for a metastable value to settle, and costs one more clock
cycle of latency on every synchronized input. At a few megahertz that extra clock period is an
enormous amount of time in metastability terms, which is exactly why the proposal sounds prudent,
and also why it may already be unnecessary.

\begin{parts}
  \item Work the trade out in both directions. Under what circumstances does it come out in favour
    of the third stage, and under what circumstances is the second stage already enough? Your
    answer has to turn on something other than the clock frequency, because frequency alone does
    not settle it: name what it does turn on.
  \item \Secref{c4:a4} got to MTBF from the resolution time available per stage. Without redoing
    that arithmetic, say which way each of these pushes the decision, and why:
    \begin{itemize}
      \item How often the asynchronous input actually changes.
      \item The metastability characteristics of the FPGA family you are targeting.
      \item What a synchronization failure would cost if one happened.
    \end{itemize}
  \item Give your recommendation for an ordinary low-speed design like the ones in this book. Then
    describe a design, real or imagined, where you would argue for the third stage instead, and say
    which of the factors in \textbf{b)} is doing the work in that case.
\end{parts}
\textbf{Check your answer covers} the two things a third stage provides, an additional clock period
of resolution time and a lower probability that an unresolved value reaches the functional logic,
and the two it costs, an additional cycle of latency and additional flip-flops. If your answer to
\textbf{a)} rests on the clock rate rather than on the required mean time between failures, read
\secref{c4:a4} again before moving on.

\exercise{Assert async, release sync}{Paper}
\label{c4:ex:5}
Explain the \textbf{assert asynchronously, release synchronously} reset pattern used in
\file{reset_sync.vhd} (\secref{c4:a5}).

\noindent Describe the two parts separately:
\begin{itemize}
  \item Asynchronous assertion:
    \begin{itemize}
      \item The reset takes effect immediately.
      \item It does not wait for a clock edge.
    \end{itemize}
  \item Synchronous release:
    \begin{itemize}
      \item The reset is removed only on an active clock edge.
      \item All affected logic leaves reset in a controlled relationship with the clock.
    \end{itemize}
\end{itemize}
Explain:
\begin{itemize}
  \item Why immediate reset assertion can be important when:
    \begin{itemize}
      \item The clock is stopped.
      \item The clock is unstable.
      \item The system must enter its reset state without delay.
    \end{itemize}
  \item Why releasing an asynchronous reset close to a clock edge can violate the recovery or
    removal timing requirements of a flip-flop.
  \item Why synchronizing the release prevents different flip-flops from leaving reset on different
    clock cycles.
\end{itemize}

% ------------------------------------------------------------------------------------------
\exerciseset{Building the Circuit}
\label{c4:set:3}

\exercise{\code{led_toggle_sync} by hand}{CircuitVerse, Paper}
\label{c4:ex:6}
Build the \code{led_toggle_sync} circuit from \secref{c4:a9} by hand in CircuitVerse.

\begin{parts}
  \item Add the external ports:
    \begin{itemize}
      \item Inputs:
        \begin{itemize}
          \item \code{clock}
          \item \code{reset_n}: asynchronous, active-low.
          \item \code{button_n}: active-low.
        \end{itemize}
      \item Output:
        \begin{itemize}
          \item \code{led}
        \end{itemize}
    \end{itemize}
  \item Build the two-flip-flop reset synchronizer:
    \begin{itemize}
      \item Connect the two flip-flops in series.
      \item Assert both stages asynchronously using the raw \code{reset_n} input.
      \item Allow the reset to propagate out of the chain synchronously.
      \item Use the synchronized reset output, \code{reset_s2_n}, throughout the rest of the
        circuit.
      \item Do not use the raw \code{reset_n} input directly outside the reset synchronizer.
    \end{itemize}
  \item Build the three-flip-flop button synchronizer and edge detector:
    \begin{itemize}
      \item Use the first two flip-flops as the double-flop synchronizer.
      \item Use the third flip-flop to store the previous synchronized button value.
      \item Compare the current and previous synchronized values.
      \item Detect a button press when:
        \begin{itemize}
          \item The current synchronized active-low button value is \code{0}.
          \item The previous synchronized value is \code{1}.
        \end{itemize}
      \item Use an AND gate with the current value inverted to generate \code{button_edge_s2}.
    \end{itemize}
    The falling-edge relationship for the active-low button is:
    \[ \bool{button\_edge\_s2} = \bool{previous} \cdot \bool{current}' \]
  \item Build the LED toggle logic:
    \begin{itemize}
      \item Add a flip-flop to store \code{led_s}.
      \item Reset \code{led_s} using \code{reset_s2_n}.
      \item Toggle \code{led_s} only when \code{button_edge_s2 = 1}.
      \item Connect the \code{led} output to \code{led_s}.
    \end{itemize}
  \item Simulate the complete design:
    \begin{itemize}
      \item Set the clock period to \code{1000 ms}.
      \item Assert and release \code{reset_n}.
      \item Press and release \code{button_n}.
      \item Observe the synchronized button stages.
      \item Observe \code{button_edge_s2}.
      \item Observe \code{led_s}.
    \end{itemize}
    Determine:
    \begin{itemize}
      \item Whether the LED toggles on the first clock edge after the physical button press.
      \item How many clock cycles are required for the button value to propagate through the
        synchronizer and edge detector.
      \item Whether the observed latency matches the behaviour predicted in \secref{c4:a9}.
    \end{itemize}
\end{parts}

% ------------------------------------------------------------------------------------------
\exerciseset{Debouncing and Timing}
\label{c4:set:4}

\exercise{Synchronized is not debounced}{Paper}
\label{c4:ex:7}
You have now built and watched the circuit, so this is a question about what you saw.

The three-flip-flop chain in Exercise~\ref{c4:ex:6} does three jobs at once:
\begin{itemize}
  \item Reduces the risk of metastability.
  \item Stores the previous synchronized button value.
  \item Detects a button press.
\end{itemize}
Explain why it does \textbf{not}, by itself, guarantee that the button is debounced.

\begin{parts}
  \item Distinguish the three things it is easy to run together:
    \begin{itemize}
      \item Synchronization: converts an asynchronous input into a signal suitable for use in the
        clock domain.
      \item Edge detection: generates a pulse when the sampled signal changes.
      \item Debouncing: prevents several physical transitions from being interpreted as several
        separate presses.
    \end{itemize}
    Which two of these does your circuit actually do?
  \item In your simulation, one press produced one clean toggle. That is a stronger claim than it
    looks: it holds only under a particular condition about where the bouncing fell relative to the
    clock edges. State that condition precisely, in terms of what your circuit sampled and what it
    therefore saw.

    Then explain why that condition is an accident of your clock period rather than a property of
    your circuit. Be specific about which of the two, the bounce or the sampling, your design
    actually controls.
  \item Your simulation ran at a \code{1000 ms} clock period; the FPGA runs at \code{50 MHz}. A push
    button's contacts typically bounce for something on the order of 1\,ms. Work out the
    consequences yourself:
    \begin{itemize}
      \item How many times does a \code{1000 ms} clock sample the button during that 1\,ms of
        bouncing?
      \item How many times does a \code{50 MHz} clock sample it?
      \item Each sampled transition looks like a fresh edge to your edge detector. Given those two
        numbers, predict what the LED does on the board for one physical press, and say why the
        identical circuit behaved impeccably in CircuitVerse.
    \end{itemize}
  \item Describe one fix of each kind:
    \begin{itemize}
      \item A hardware fix: for example, an RC low-pass filter followed by a Schmitt-trigger input.
      \item A purely digital fix:
        \begin{itemize}
          \item for example, requiring the synchronized signal to hold a new level for some minimum
            number of clock cycles before accepting it.
          \item or a counter that ignores further transitions for a fixed interval after a press,
            which is what Chapter~\ref{c7:ch}'s timer would give you.
        \end{itemize}
    \end{itemize}
  \item Explain why the double-flop synchronizer must still be kept even once a debounce method is
    added. What problem does debouncing not solve?
\end{parts}

% ------------------------------------------------------------------------------------------
\exerciseset{The Synchronizers as Modules of Their Own}
\label{c4:set:5}
The worked example puts each synchronizer in a module of its own, so every later design
instantiates the same two files instead of re-deriving them. That is why Chapter~\ref{c7:ch}'s
\code{walking_led} and Chapter~\ref{c8:ch}'s \code{fsm_led} both ask you to copy these two in, and
why Chapter~\ref{c8:ch}'s \code{seq_detect_mealy}, which has no button, takes the reset synchronizer
alone and leaves the other behind (\secref{c4:a9}).

The next two exercises are you writing those modules. \Secref{c4:a5} and \secref{c4:a6} give you
the logic; what they do not give you is the entity around it, and, for the second one, the step
from a single button to a vector of them. Write each module yourself before opening the worked
example's copy, then compare. Exercise~\ref{c4:ex:10} then composes both of them into a working
design.

\exercise{The reset synchronizer}{VHDL, Paper}
\label{c4:ex:8}
Write the reset synchronizer as a module named \code{reset_sync}.

\noindent The entity has:

\begin{booktable}{@{}p{6em}p{4.9em}p{5.2em}L@{}}
  \thead{Port} & \thead{Direction} & \thead{Type} & \thead{Description} \\
  \midrule
  \code{clock} & in & \code{std_logic} & 50\,MHz system clock. \\
  \code{reset_n} & in & \code{std_logic} & Active-low asynchronous reset, straight from the
    outside world. \\
  \code{reset_s2_n} & out & \code{std_logic} & Active-low synchronized reset, safe for the rest of
    the design to use. Goes low the instant \code{reset_n} does, with no clock edge involved, and
    returns high two rising edges after \code{reset_n} does. \\
\end{booktable}

\begin{parts}
  \item Implement \textbf{assert asynchronously, release synchronously} (\secref{c4:a5}):
    \begin{itemize}
      \item One internal signal for the first stage, plus the output for the second.
      \item One process, sensitive to both \code{clock} and \code{reset_n}.
      \item On \code{reset_n = '0'}, drive both stages low, outside the clocked branch.
      \item On a rising clock edge, shift a constant \code{'1'} in through the two stages.
    \end{itemize}
  \item The value forced on reset and the value shifted in on the clock are opposites, \code{'0'}
    and \code{'1'}. Explain in one sentence why that is not a contradiction, in terms of what each
    branch is for.
  \item Verify the design with its testbench (see the aside at the start of this section). It checks
    both halves of the pattern separately: that pulling \code{reset_n} low takes \code{reset_s2_n}
    low with no clock edge in between, and that releasing \code{reset_n} leaves \code{reset_s2_n}
    low until two rising edges have passed.
  \item Answer, in prose:
    \begin{itemize}
      \item The output is named \code{reset_s2_n} rather than \code{reset_n_sync}. What does
        \code{s2} say, and what would have to change about this module for that name to become
        wrong?
      \item Every other module in this book takes \code{reset_s2_n} as its reset. This one takes
        the raw \code{reset_n}. Why must exactly one module in a design be the exception?
    \end{itemize}
\end{parts}
\note[Tip] the release is the interesting half. On reset the process does not care about the clock
at all, so both stages go low together; on the clock it does not care about \code{reset_n}, so the
\code{'1'} takes one edge per stage to reach the output. Two edges of latency on release is the
price of the whole pattern.

\modulefigure{reset_sync}

\selfcheck name your entity \code{reset_sync}, with ports \code{clock}, \code{reset_n} (in) and
\code{reset_s2_n} (out), declared in that order; its testbench is in \file{exercises/reset_sync}.

\exercise{The button synchronizer}{VHDL, Paper}
\label{c4:ex:9}
Write the button synchronizer and edge detector as a module named \code{button_sync}.

Exercise~\ref{c4:ex:10} handles two buttons by writing a two-bit synchronizer. This one is the same
circuit with the width left open: one generic decides how many buttons it serves, so
\code{led_toggle_sync} can use it for one button and \code{fsm_led} for two without either file
being edited. See \secref{c4:a8} for the syntax and for why this module has a generic where
\code{reset_sync} has none.

\noindent The entity has one generic:

\begin{booktable}{@{}p{4.5em}p{10.4em}p{3.8em}L@{}}
  \thead{Generic} & \thead{Type} & \thead{Default} & \thead{Description} \\
  \midrule
  \code{COUNT} & \code{natural range 1 to 3} & \code{1} & The number of buttons, and therefore the
    width of both vector ports. \\
\end{booktable}

\noindent and these ports:

\begin{booktable}{@{}p{7.4em}p{4.9em}>{\raggedright\arraybackslash}p{11.4em}L@{}}
  \thead{Port} & \thead{Direction} & \thead{Type} & \thead{Description} \\
  \midrule
  \code{clock} & in & \code{std_logic} & 50\,MHz system clock. \\
  \code{reset_s2_n} & in & \code{std_logic} & Active-low, \textbf{already synchronized} reset.
    This module does not synchronize its own reset; Exercise~\ref{c4:ex:8}'s module did that. \\
  \code{button_n} & in & \code{std_logic_vector(COUNT - 1 downto 0)} & Active-low asynchronous
    push buttons. \\
  \code{button_edge_s2} & out & \code{std_logic_vector(COUNT - 1 downto 0)} & A one-cycle pulse
    on each button's falling edge, two clock edges after the press. Each bit is independent of
    every other. \\
\end{booktable}

\begin{parts}
  \item Build the three-stage chain from \secref{c4:a6}, now over vectors rather than single bits:
    \begin{itemize}
      \item Three internal \code{std_logic_vector(COUNT - 1 downto 0)} signals, one per stage.
      \item Stages 1 and 2 are the double-flop synchronizer; stage 3 stores the previous value of
        stage 2.
      \item On reset, set all three stages to \code{(others => '1')}, not \code{(others => '0')}.
      \item Derive the output with a single concurrent assignment, and note that the same
        expression works unchanged on vectors as on bits.
    \end{itemize}
    \textbf{Keep this one in mind.} Those first two flip-flops are the whole of a synchronizer and
    nothing else, and Chapter~\ref{c5:ch}'s Exercise~\ref{c5:ex:4} has you pull them back out as a
    module of their own, generic in width and in the value they reset to. Writing them inline here
    first is what makes that exercise a consolidation rather than a new idea.
  \item Explain why the reset value is all-ones. All-ones means ``released'', because the buttons
    are active-low, so an all-zeros reset starts the chain out claiming every button is already
    held down. Trace what that costs:
    \begin{itemize}
      \item With the buttons released and a reset to all-zeros, does the module emit a spurious
        pulse as the released state propagates through the three stages? Work it through the
        output expression edge by edge before you answer, rather than guessing.
      \item Now the case that actually breaks. A user is holding a button down at the moment the
        reset is released. What does the module report with an all-ones reset, and what does it
        report with all-zeros? Which of the two is the behaviour you want, and why?
    \end{itemize}
  \item Verify the design with its testbench. It instantiates your module \textbf{twice}, once at
    \mbox{\code{COUNT = 1}} and once at \mbox{\code{COUNT = 2}}, and checks that a press produces
    exactly one single-cycle pulse, that holding the button produces no further pulses, that a
    release produces none at all, that pressing one button of a pair never pulses the other, and
    --- the case part \textbf{b)} asks you to reason about --- that a button already held down when
    the reset releases produces exactly one pulse with an all-ones reset and none at all with
    all-zeros.
  \item Answer, in prose:
    \begin{itemize}
      \item \code{COUNT} is constrained to \code{1 to 3} rather than left an unbounded
        \code{natural}. What does the constraint buy, given that nothing in the architecture
        depends on the upper bound?
      \item This module is deliberately not merged with Exercise~\ref{c4:ex:8}'s, even though a
        design that has a button always needs both. Give the design that justifies the split, and
        say what it would have had to do instead.
    \end{itemize}
\end{parts}
\note[Tip] write the architecture for \code{COUNT = 1} in your head first, then check that every
line you wrote is already correct for a vector. \code{and}, \code{not} and the aggregate
\code{(others => '1')} all work element-wise, so the generalization should cost you no extra code at
all. If it does, that is worth a second look.

\modulefigure{button_sync}

\selfcheck name your entity \code{button_sync}, with generic \code{COUNT}
(\code{natural range 1 to 3}, default \code{1}), inputs \code{clock}, \code{reset_s2_n},
\code{button_n} (\code{std_logic_vector(COUNT - 1 downto 0)}) and output \code{button_edge_s2}
(\code{std_logic_vector(COUNT - 1 downto 0)}), declared in that order; its testbench is in
\file{exercises/button_sync}.

% ------------------------------------------------------------------------------------------
\exerciseset{Composing Them Into a Design}
\label{c4:set:6}

\exercise{Two synchronized LED toggles}{VHDL, Paper}
\label{c4:ex:10}
Now make Chapter~\ref{c3:ch}'s \code{led_toggle} metastability-safe, in VHDL, as a new module named
\code{led_toggle_sync2}.

Chapter~\ref{c3:ch}'s \code{led_toggle} toggled two LEDs from two push buttons, but fed their raw,
asynchronous signals straight into the edge detector: logically correct, but unsafe on real
hardware (Chapter~\ref{c3:ch}'s own ``Looking ahead'' flagged exactly this). Here you rebuild it
with every asynchronous input guarded by this chapter's double-flop synchronizer.

\noindent The entity has:

\begin{booktable}{@{}p{5.2em}p{4.9em}p{14.4em}L@{}}
  \thead{Port} & \thead{Direction} & \thead{Type} & \thead{Description} \\
  \midrule
  \code{clock} & in & \code{std_logic} & 50\,MHz system clock. \\
  \code{reset_n} & in & \code{std_logic} & Active-low asynchronous reset. \\
  \code{button_n} & in & \code{std_logic_vector(1 downto 0)} & Two active-low asynchronous push
    buttons. \\
  \code{led} & out & \code{std_logic_vector(1 downto 0)} & Two LEDs; \code{led(i)} toggles once per
    press of \code{button_n(i)}. \\
\end{booktable}

\begin{parts}
  \item Synchronize every asynchronous input:
    \begin{itemize}
      \item Apply the two-flip-flop reset synchronizer from \secref{c4:a5} (assert asynchronously,
        release synchronously).
      \item Give the buttons a double-flop synchronizer plus a third stage for edge detection, as
        in \secref{c4:a6}, now over a \code{std_logic_vector(1 downto 0)}, so both buttons are
        synchronized in parallel.
      \item Never let the raw \code{reset_n} or \code{button_n} reach the toggle logic; use only
        the synchronized reset and the per-button edge pulses.
    \end{itemize}
    \note[Tip] you can inline the synchronizer, or split it out into a subcomponent taking a two-bit
    \code{button_n}, using the instantiation syntax from \secref{c2:a6}. Inlining is the shorter
    route, and then the build is a single file:
\begin{shell}
ghdl -a --std=93 led_toggle_sync2.vhd led_toggle_sync2_tb.vhd
\end{shell}
    If you split it out instead, copy in the \file{reset_sync.vhd} and \file{button_sync.vhd} you
    wrote for Exercises~\ref{c4:ex:8} and~\ref{c4:ex:9} (instantiating \code{button_sync} at
    \code{COUNT = 2}), and name them first on the \code{ghdl -a} line, ahead of your own module, so
    GHDL sees them before the file that instantiates them:
\begin{shell}
ghdl -a --std=93 reset_sync.vhd button_sync.vhd \
                 led_toggle_sync2.vhd led_toggle_sync2_tb.vhd
\end{shell}
    See \secref{c2:c5}. Neither module is shipped anywhere in the repository, because writing them
    is Exercises~\ref{c4:ex:8} and~\ref{c4:ex:9}. If yours are not working yet, go back to
    \secref{c4:a5}, which prints \code{reset_sync}'s process in full, and \secref{c4:a6}, which
    gives \code{button_sync}'s three-flop chain and its edge equation.
  \item Build the toggle logic:
    \begin{itemize}
      \item One LED-state flip-flop per LED.
      \item On the synchronized reset, clear both LEDs.
      \item Toggle \code{led(i)} only when button \code{i}'s synchronized falling-edge pulse fires.
      \item Confirm \code{led(i)} responds to \code{button_n(i)} alone, never to the other button.
    \end{itemize}
  \item Verify the design with its testbench (see the aside at the start of this section). The
    testbench presses each button in turn and checks that:
    \begin{itemize}
      \item Reset clears both LEDs.
      \item No press means no toggle.
      \item Pressing button 0 toggles \code{led(0)} \textbf{only}, once the synchronizer latency has
        elapsed.
      \item Holding that button down does not toggle it again.
      \item Pressing button 1 then toggles \code{led(1)} only, leaving \code{led(0)} where it was.
    \end{itemize}
    Pay attention to the \textbf{latency}: the testbench checks the \emph{exact} edge, not merely
    that the LED settles eventually. A press must reach the toggle logic through two synchronizer
    flip-flops, so \code{led} is still unchanged after the first and second rising edges and toggles
    on the third. Count the flip-flops a press travels through and satisfy yourself that the third
    edge is the right one, and that a defensive extra pipeline stage would now be caught rather than
    tolerated.
  \item Answer, in prose, by comparing this design with your Chapter~\ref{c3:ch} \code{led_toggle}:
    \begin{itemize}
      \item What does the synchronizer change about pressing a real, bouncing button?
      \item Why was the raw Chapter~\ref{c3:ch} version unsafe to wire to a physical button at
        \code{50 MHz}?
      \item Your testbench drives clean, perfectly-timed transitions. Which of the two problems
        from this chapter, metastability or contact bounce, can a testbench therefore never
        demonstrate to you?
    \end{itemize}
\end{parts}

\modulefigure{led_toggle_sync2}

\selfcheck name your entity \code{led_toggle_sync2}, with inputs \code{clock}, \code{reset_n},
\code{button_n} (\code{std_logic_vector(1 downto 0)}) and output \code{led}
(\code{std_logic_vector(1 downto 0)}), declared in that order; its testbench is in
\file{exercises/led_toggle_sync2}. It presses each button in turn and checks that \code{led(i)}
toggles on \code{button_n(i)}'s falling edge, that it does so on the right clock edge rather than
merely settling to the right value eventually, and that it never fires the wrong LED; see
\secref{c2:app:c}.

